\chapter{Determinan dan Sistem Linear}\label{ch:b1:det}

Determinan memampatkan ke dalam satu skalar jawaban atas ``apakah
$n$ vektor ini sebuah \omterm{def:b1:vspaces:free}{basis}?'' --- dan, secara geometri, mengukur volume
yang direntangnya. Kita mencirikannya lewat sifatnya (multilinear,
berselang-seling, ternormalkan), menghitungnya dalam dimensi $2$ dan $3$ serta lewat
\omterm{thm:b1:det:cofactor}{penjabaran kofaktor} pada umumnya, lalu menerapkannya pada \omterm{def:b1:det:system}{sistem linear},
berdampingan dengan algoritma serbaguna: yaitu \omterm{met:b1:det:gauss}{penghapusan Gauss}.

\section{Determinan}

\begin{theorem}[Pencirian]\label{thm:b1:det:def}
Ada tepat satu \omterm{def:b1:logic:map}{pemetaan} $\det \colon \mathcal{M}_n(K) \to K$,
yang dipandang sebagai fungsi atas $n$ kolomnya, yang:
\begin{enumerate}
  \item \emph{\omterm{def:b1:linmaps:def}{linear} pada setiap kolomnya} (dengan yang lain tetap);
  \item \emph{berselang-seling}: menukar dua kolom mengubah tandanya
        (sehingga dua kolom yang sama memberikan $0$);
  \item \emph{ternormalkan}: $\det I_n = 1$.
\end{enumerate}
Untuk $n = 2$ dan $3$:\index{determinan}
\[
\begin{vmatrix} a & b\\ c & d\end{vmatrix} = ad - bc,
\qquad
\begin{vmatrix} a & b & c\\ d & e & f\\ g & h & i\end{vmatrix}
= aei + bfg + cdh - ceg - bdi - afh
\]
(yaitu kaidah Sarrus $3 \times 3$: hasil kali diagonal yang menurun
dikurangi yang menaik).
\end{theorem}
\admitted

\begin{remark}
Untuk $n = 2$: menjabarkannya lewat kebilinearan pada kolom kanoniknya memberikan
rumusnya, yang sebaliknya memenuhi aksiomanya --- jadi bukti yang
lengkap; sedangkan $n = 3$ serupa dengan lebih banyak suku. Adapun kasus umumnya
(dengan keberadaannya lewat jumlah atas \omterm{def:b1:counting:objects}{permutasi}, dan ketunggalannya lewat penjabaran
yang sama) menuntut tanda sebuah \omterm{def:b1:counting:objects}{permutasi} dan ditangguhkan ke
tahun kedua; jadi kita \omterm{def:b1:vspaces:free}{bebas} memakai aksioma dan akibatnya
di bawah.

Berikut penjabaran $n = 2$ selengkapnya, karena ia menjadi contohnya: dengan
kolom $C_1 = a\,e_1 + c\,e_2$ dan $C_2 = b\,e_1 + d\,e_2$,
kebilinearannya memberikan
\[
\det(C_1, C_2) = ab\det(e_1, e_1) + ad\det(e_1, e_2)
+ cb\det(e_2, e_1) + cd\det(e_2, e_2),
\]
dan sifat berselang-selingnya membunuh pasangan yang berulang sembari membalik
$\det(e_2, e_1) = -\det(e_1, e_2)$: sehingga seluruh \omterm{def:b1:logic:map}{pemetaannya} runtuh menjadi
$(ad - bc)\det(e_1, e_2) = ad - bc$ lewat penormalannya. Adapun ketunggalannya
kasatmata pada perhitungannya sendiri --- karena aksiomanya tak menyisakan
pilihan pada langkah mana pun --- dan inilah persis fakta ketunggalan
berskala yang dipakai pada bukti kaidah hasil kali di bawah.
\end{remark}

\begin{theorem}[Sifat]\label{thm:b1:det:props}
Untuk $A, B \in \mathcal{M}_n(K)$:
\begin{enumerate}
  \item menambahkan kelipatan sebuah kolom pada kolom lain tak mengubah
        determinannya; sedangkan mengalikan sebuah kolom dengan $\lambda$ mengalikannya
        dengan $\lambda$ (sehingga $\det(\lambda A) = \lambda^n \det A$);
  \item $\det(AB) = \det A\, \det B$;
  \item $A$ dapat dibalik $\iff$ $\det A \neq 0$ $\iff$ kolomnya
        membentuk \omterm{def:b1:vspaces:free}{basis} $K^n$; dan lalu $\det(A^{-1}) = (\det
        A)^{-1}$;
  \item $\det(A^{\mathsf T}) = \det A$ --- sehingga setiap kaidah kolomnya
        juga kaidah baris;
  \item determinan sebuah matriks segitiga adalah hasil kali entri
        diagonalnya.
\end{enumerate}
\end{theorem}

\begin{proof}
(1) Lewat kelinearannya, $\det(\dots, C_i + \lambda C_j, \dots) = \det A +
\lambda\det(\dots, C_j, \dots)$ dan determinan keduanya berkolom dua
yang sama: jadi nol.

(2) Tetapkanlah $A$ lalu pandanglah $\varphi(B) = \det(AB)$ sebagai fungsi atas
kolom $B$: karena $AB$ berkolom $AB_j$, maka $\varphi$ bersifat
multilinear dan berselang-seling pada $B_j$. Kita mengakui, bersama
\cref{thm:b1:det:def}, \omterm{def:b1:logic:statement}{pernyataan} ketunggalannya dalam bentuk berskala:
bahwa \emph{setiap} \omterm{def:b1:logic:map}{pemetaan} multilinear berselang-seling $\varphi$ atas kolomnya
sama dengan $\varphi(I_n) \cdot \det$. Di sini $\varphi(I_n) = \det A$, sehingga
$\det(AB) = \det A \cdot \det B$.

(3) Jika $A$ dapat dibalik: maka $\det A\,\det A^{-1} = \det I = 1$, sehingga
$\det A \neq 0$ dan rumus inversnya berlaku. Jika $A$ tak
dapat dibalik, maka kolomnya terkait (\cref{cor:b1:linmaps:samedim}
dan \cref{prop:b1:linmaps:basis}); lalu menyatakan satu kolom lewat
yang lain dan menjabarkannya secara \omterm{def:b1:linmaps:def}{linear} menyisakan determinan berkolom dua
yang sama: jadi $\det A = 0$. Adapun \omterm{def:b1:logic:statement}{pernyataan} basisnya adalah
\cref{prop:b1:findim:twoofthree}.

(4) Diakui bersama konstruksi umumnya (karena ia seketika pada
rumus \omterm{def:b1:counting:objects}{permutasinya}); kita mencatatnya agar dapat memakai \omterm{met:b1:matrices:gauss}{operasi baris}.

(5) Jika suatu entri diagonalnya lenyap, maka $k$ kolom pertamanya
terkait untuk suatu $k$ (lewat pertimbangan rank) dan $\det = 0 =$ hasil kalinya.
Kalau tidak, bersihkanlah setiap kolomnya di bawah-dan-kiri lewat operasi jenis
(1) --- yang mungkin dalam bentuk segitiganya --- sampai mencapai matriks
diagonal, yang determinannya adalah hasil kali entrinya lewat
kemultilinearan dari $I_n$.
\end{proof}

\begin{example}[Kaidahnya, diperiksa pada bilangan]
\label{ex:b1:det:rulescheck}
Ambillah $A = \begin{pmatrix} 1 & 2\\ 3 & 4\end{pmatrix}$ ($\det A =
-2$) dan $B = \begin{pmatrix} 0 & 1\\ 1 & 1\end{pmatrix}$ ($\det
B = -1$). Maka
\[
AB = \begin{pmatrix} 2 & 3\\ 4 & 7\end{pmatrix},
\quad \det(AB) = 14 - 12 = 2 = (-2)(-1) ;
\qquad
\det(A^{\mathsf T}) = \begin{vmatrix} 1 & 3\\ 2 & 4
\end{vmatrix} = -2 = \det A .
\]
Kesifatan perkalian dan keinvarianan transposnya tertegaskan --- sedangkan
kesifatan penjumlahan yang \emph{palsu} gagal pada pasangan yang sama:
\[
\det(A + B) = \begin{vmatrix} 1 & 3\\ 4 & 5\end{vmatrix} = -7
\neq \det A + \det B = -3 .
\]
Tiga puluh detik aritmetika semacam ini, setelah sebarang kesamaan
determinan dipanggil, merupakan asuransi galat termurah yang ada.
\end{example}

\begin{example}[Determinan sebagai luas]
\label{ex:b1:det:area}
Jajargenjang yang direntang $u = (2, 0)$ dan $v = (1, 3)$ mempunyai
alas $2$ dan tinggi $3$: jadi luasnya $6$. Dan
\[
\begin{vmatrix} 2 & 1\\ 0 & 3\end{vmatrix} = 6 :
\]
determinan $2\times2$ itu \emph{adalah} luas bertanda pada
jajargenjang kolomnya. Adapun aksiomanya menuturkan ulang geometrinya:
menambahkan kelipatan satu kolom pada yang lain merupakan \emph{geseran},
yang menggelincirkan jajargenjangnya sejajar sebuah sisi tanpa mengubah
alas atau tingginya (yaitu operasi~(1) pada \cref{thm:b1:det:props});
menskalakan sebuah kolom menskalakan luasnya; sedangkan menukar kolomnya membalik
arahnya, dan dari situlah tandanya, $\det(v, u) = -6$. Di dalam $\R^3$
pembacaan yang sama memberikan volume bertanda, dan $\abs{\det}$ menjadi
faktor penskalaan volume yang universal bagi \omterm{def:b1:linmaps:def}{pemetaan linear} --- yaitu fakta
di balik rumus perubahan variabel bagi \omterm{thm:b1:integration:def}{integral} rangkap
pada jilid Tahun~2.
\end{example}

\begin{theorem}[Penjabaran kofaktor]\label{thm:b1:det:cofactor}
Misalkan $A \in \mathcal{M}_n(K)$ dan $\Delta_{ij}$ determinan
$A$ dengan baris $i$ dan kolom $j$ yang dihapus. Maka, untuk sebarang kolom tetap
$j$ (atau baris, lewat transposisi):\index{penjabaran kofaktor}
\[
\det A = \sum_{i=1}^{n} (-1)^{i+j}\, a_{ij}\, \Delta_{ij} .
\]
\end{theorem}
\admitted

\begin{example}\label{ex:b1:det:cofactor}
Menjabarkannya sepanjang kolom pertamanya:
\[
\begin{vmatrix}
2 & 1 & 0\\
1 & 2 & 1\\
0 & 1 & 2
\end{vmatrix}
= 2\begin{vmatrix} 2 & 1\\ 1 & 2\end{vmatrix}
- 1\begin{vmatrix} 1 & 0\\ 1 & 2\end{vmatrix}
= 2 \times 3 - 2 = 4 .
\]
Strateginya: ciptakanlah nol dahulu (lewat \omterm{met:b1:matrices:gauss}{operasi baris}/kolom), lalu jabarkanlah
sepanjang garis yang paling kosong.
\end{example}

\begin{example}[Invers kofaktornya, sekali dengan tangan]
\label{ex:b1:det:cofactorinverse}
Untuk $A = \begin{pmatrix} 1 & 1 & 0\\ 0 & 1 & 1\\ 1 & 0 &
1\end{pmatrix}$: $\det A = 1(1) - 1(-1) + 0 = 2$. Adapun kesembilan
kofaktornya $(-1)^{i+j}\Delta_{ij}$ terhimpun menjadi
\[
\operatorname{Com}(A) = \begin{pmatrix}
1 & 1 & -1\\
-1 & 1 & 1\\
1 & -1 & 1
\end{pmatrix},
\qquad
A^{-1} = \frac{1}{\det A}\operatorname{Com}(A)^{\mathsf T}
= \frac12\begin{pmatrix}
1 & -1 & 1\\
1 & 1 & -1\\
-1 & 1 & 1
\end{pmatrix},
\]
yaitu rumus yang dikutip pada \cref{exo:b1:det:8}. Periksalah satu pasangan
baris-kolom: (baris $1$ pada $A$)(kolom $1$ pada $A^{-1}$) $= \frac12(1 + 1
+ 0) = 1$, dan terhadap kolom $2$: $\frac12(-1 + 1 + 0) = 0$.
Sembilan determinan $2\times2$ bagi satu invers $3\times3$: sehingga sudah
pada ukuran ini, reduksi baris (\cref{exo:b1:det:3}) lebih murah ---
adapun nilai rumus kofaktornya bersifat teori (yaitu kebulatan pada
\cref{exo:b1:det:8}, dan sifat \omterm{def:b1:derivative:def}{dapat diturunkan} inversnya pada jilid
kemudian), bukan perhitungan.
\end{example}

\begin{example}[Kaidah blok segitiganya, pada ukuran $4$]
\label{ex:b1:det:blocktriangular}
Klaimnya: $\det\begin{pmatrix} M & N\\ 0 & P\end{pmatrix} = \det
M\,\det P$ bagi blok $2\times2$. Bersihkanlah blok $N$ lewat operasi
kolom: karena menambahkan pada kolom $3, 4$ kombinasi kolom
$1, 2$ yang cocok membuang $N$ \emph{ketika $M$ dapat dibalik} (pecahkanlah
$M\Lambda = -N$ bagi koefisien kombinasinya $\Lambda$),
sehingga menyisakan $\det\begin{pmatrix} M & 0\\ 0 & P\end{pmatrix}$; lalu
\omterm{thm:b1:det:cofactor}{penjabaran kofaktor} sepanjang kolom pertamanya, dua kali, memberikan $\det
M\det P$ bagi bentuk blok diagonal ini. Jika $M$ tak
dapat dibalik, maka kolomnya terkait, sehingga kedua kolom pertama
matriks besarnya terkait (karena separuh bawahnya nol): jadi kedua
sisinya lenyap. Kaidahnya meluas ke sebarang ukuran blok lewat argumen
dua kasus yang sama --- dan ialah mesin
\cref{exo:b1:det:10}.
\end{example}

\begin{example}[Sebuah determinan $4 \times 4$, secara strategis]
\label{ex:b1:det:fourbyfour}
\[
\Delta = \begin{vmatrix}
1 & 2 & 3 & 4\\
2 & 3 & 4 & 1\\
3 & 4 & 1 & 2\\
4 & 1 & 2 & 3
\end{vmatrix}.
\]
Setiap barisnya berjumlah $10$: sehingga operasi $C_1 \leftarrow C_1 + C_2 +
C_3 + C_4$ membuat kolom pertamanya tetap, dan memfaktorkan $10$
menyisakan satuan. Lalu $L_i \leftarrow L_i - L_1$ ($i \geq 2$) membersihkan
kolom pertamanya:
\[
\Delta = 10\begin{vmatrix}
1 & 2 & 3 & 4\\
0 & 1 & 1 & -3\\
0 & 2 & -2 & -2\\
0 & -1 & -1 & -1
\end{vmatrix}
= 10\begin{vmatrix}
1 & 1 & -3\\
2 & -2 & -2\\
-1 & -1 & -1
\end{vmatrix}
= 10 \times 16 = 160,
\]
dengan determinan $3\times3$ terakhirnya dijabarkan sepanjang baris pertamanya:
$1(2 - 2) - 1(-2 - 2) + (-3)(-2 - 2) = 0 + 4 + 12 = 16$. Moralnya:
satu operasi yang terpilih baik (yaitu melihat jumlah baris yang tetap) mengalahkan
enam belas kofaktor.
\end{example}

\begin{method}[Memilih strategi determinan]
\label{met:b1:det:strategy}
Pindailah matriksnya sebelum menghitung apa pun.
\begin{enumerate}
  \item \emph{Jumlah baris atau kolom yang tetap}: tambahkanlah segalanya ke dalam
        satu garis, lalu faktorkanlah nilai bersamanya
        (\cref{ex:b1:det:fourbyfour}, \cref{exo:b1:det:7}).
  \item \emph{Struktur yang berulang}: kurangkanlah baris atau kolom yang
        bertetangga untuk menciptakan nol; karena pola tangganya runtuh
        menuju bentuk segitiga, yang determinannya terbaca pada
        diagonalnya.
  \item \emph{Nol yang terpencil}: jabarkanlah sepanjang garis yang paling kosong
        (\cref{ex:b1:det:cofactor}); sedangkan keluarga yang berulang
        (yaitu tridiagonal, \cref{exo:b1:det:6}) menghasilkan rekurensi lewat
        cara ini.
  \item \emph{Sebuah parameter}: determinannya berupa \omterm{def:b1:poly:def}{polinomial} pada
        parameter itu; carilah akarnya dengan melihat nilai yang merosot
        (yaitu baris yang sama, kolom yang sebanding), lalu pakukanlah
        \omterm{def:b1:poly:def}{polinomialnya} lewat derajat dan koefisien utamanya. Bagi
        matriks pada \cref{exo:b1:det:7}: $m = 1$ memberikan tiga
        baris yang sama (ber-rank $1$, yaitu akar rangkap), sedangkan $m = -2$ membuat
        jumlah barisnya nol (jadi satu akar lagi); adapun determinannya
        berderajat $3$ pada $m$ dengan suku utama $-m^3$ (yaitu
        hasil kali antidiagonalnya $m\cdot m\cdot m$, yang bertanda
        Sarrus $-1$), sehingga ia pastilah $-(m+2)(m-1)^2$ --- jadi tak ada
        penjabaran yang diperlukan, dan kedua metodenya saling memeriksa.
\end{enumerate}
\end{method}

\begin{example}[Determinan Vandermonde]\label{ex:b1:det:vandermonde}
Untuk skalar $x_1, \dots, x_n$:\index{determinan Vandermonde}
\[
V(x_1, \dots, x_n) =
\begin{vmatrix}
1 & x_1 & x_1^2 & \cdots & x_1^{n-1}\\
1 & x_2 & x_2^2 & \cdots & x_2^{n-1}\\
\vdots & & & & \vdots\\
1 & x_n & x_n^2 & \cdots & x_n^{n-1}
\end{vmatrix}
= \prod_{1 \leq i < j \leq n} (x_j - x_i) .
\]
Sketsa buktinya (yang terperinci pada \cref{exo:b1:det:5}): operasi kolom
$C_k \leftarrow C_k - x_1 C_{k-1}$ dari kanan membersihkan baris
pertamanya, dan memfaktorkan setiap baris sisanya mereduksinya menjadi
$V(x_2, \dots, x_n)$. Ia taknol jika dan hanya jika $x_i$ berbeda berpasangan
--- yaitu determinan di balik \omterm{thm:b1:poly:lagrange}{interpolasi Lagrange}
(\cref{ex:b1:linmaps:interpolation}).
\end{example}

\section{Sistem linear}

\begin{definition}\label{def:b1:det:system}
Sebuah sistem \omterm{def:b1:linmaps:def}{linear}\index{sistem linear} berisi $n$ persamaan dengan $p$
variabel adalah $AX = B$ dengan $A \in \mathcal{M}_{n,p}(K)$, $B \in
K^n$; ia disebut \emph{homogen} ketika $B = 0$. Adapun \omterm{def:b1:logic:sets}{himpunan} penyelesaiannya, ketika
tak kosong, adalah $X_0 + \ker A$: yaitu penyelesaian khusus tambah penyelesaian
homogen umumnya --- jadi \omterm{def:b1:vspaces:subspace}{subruang} afin berdimensi $p -
\operatorname{rk} A$ (lewat rank--nulitas).
\end{definition}

\begin{example}[Struktur afinnya, dibuat kasatmata]
\label{ex:b1:det:affinestructure}
Pecahkanlah
\[
\begin{cases}
x + y + z = 3\\
x - y + 2z = 2 .
\end{cases}
\]
Dengan mengurangkan persamaannya: $2y - z = 1$, sehingga $z = 2y - 1$ dan $x =
3 - y - z = 4 - 3y$. Maka penyelesaiannya membentuk garis
\[
(x, y, z) = (4,\ 0,\ -1) + y\,(-3,\ 1,\ 2)
\qquad (y \in \R):
\]
yaitu penyelesaian khususnya $X_0 = (4, 0, -1)$ (dari pilihan $y = 0$)
tambah garis kernelnya $\ker A = \operatorname{Vect}(-3, 1, 2)$ pada
sistem homogen yang berkaitan --- periksa: $(-3) + 1 + 2 = 0$
dan $(-3) - 1 + 4 = 0$. Secara geometri, dua bidang yang tak sejajar pada
$\R^3$ beririsan sepanjang sebuah garis, dan cacah dimensinya $p -
\operatorname{rk} A = 3 - 2 = 1$ sudah tahu itu sebelum kita memecahkan
apa pun. Mengubah penyelesaian khususnya (katakanlah $y = 1$: $X_0' =
(1, 1, 1)$) mengubah pemeriannya, bukan garisnya: karena \omterm{def:b1:vspaces:subspace}{subruang}
afin berasal-usul banyak dan berarah satu.
\end{example}

\begin{theorem}[Sistem Cramer persegi]\label{thm:b1:det:cramer}
Jika $A \in GL_n(K)$, maka sistem $AX = B$ berpenyelesaian tunggal $X =
A^{-1}B$, yang \omterm{prop:b1:vspaces:coordinates}{koordinatnya}\index{kaidah Cramer}
\[
x_j = \frac{\det A_j}{\det A},
\qquad A_j = A \text{ dengan kolom } j \text{ diganti oleh } B .
\]
\end{theorem}

\begin{proof}
Ketunggalan dan keberadaannya adalah keterbalikannya. Adapun untuk rumusnya:
tulislah $B = \sum_k x_k C_k$ (dengan kolom $A$); lalu, lewat
kemultilinearan dan sifat berselang-selingnya,
\[
\det A_j = \det\Bigl(C_1, \dots, \sum_k x_k C_k, \dots, C_n\Bigr)
= \sum_k x_k \det(C_1, \dots, C_k, \dots, C_n)
= x_j \det A ,
\]
karena setiap sukunya kecuali $k = j$ berkolom berulang.
\end{proof}

\begin{example}[Cramer dengan sebuah parameter, selengkapnya]
\label{ex:b1:det:cramerparameter}
Untuk $m \in \R$, pecahkanlah
\[
\begin{cases}
x + m y = 1\\
m x + y = 2 .
\end{cases}
\]
Determinannya adalah $1 - m^2$. \emph{Kasus umumnya} $m \neq \pm1$:
Cramer memberikan
\[
x = \frac{\begin{vmatrix} 1 & m\\ 2 & 1\end{vmatrix}}{1 - m^2}
= \frac{1 - 2m}{1 - m^2},
\qquad
y = \frac{\begin{vmatrix} 1 & 1\\ m & 2\end{vmatrix}}{1 - m^2}
= \frac{2 - m}{1 - m^2},
\]
yaitu satu penyelesaian yang bersih bagi setiap $m$ yang diperkenankan (periksa di $m = 0$:
$(1, 2)$, yang jelas benar). \emph{Kasus yang merosot}: di $m = 1$
persamaannya terbaca $x + y = 1$ dan $x + y = 2$: jadi tak sejalan; di
$m = -1$ keduanya terbaca $x - y = 1$ dan $-x + y = 2$, yaitu $x - y =
1$ dan $x - y = -2$: tak sejalan lagi. Adapun lenyapnya
determinannya mengumumkan bahwa \emph{sesuatu} merosot, tetapi
tak pernah mengatakan apa --- karena kosong atau takhingga mesti diputuskan dengan melihat
ruas kanannya. Perhatikanlah pula bagaimana rumusnya mengisyaratkan
batasnya sendiri: ketika $m \to 1^{-}$, $x = \frac{1 - 2m}{1 - m^2} \to
-\infty$; sehingga titik penyelesaiannya melarikan diri saat kedua garisnya menjadi
sejajar.
\end{example}

\begin{method}[Penghapusan Gauss pada sistem]\label{met:b1:det:gauss}
Reduksikanlah baris matriks yang diperbesar $(A \mid B)$ ke bentuk
eselon.\index{penghapusan Gauss}
\begin{enumerate}
  \item Jika sebuah poros muncul pada kolom terakhirnya (yaitu baris $0 = 1$): maka tak ada
        penyelesaian.
  \item Kalau tidak, variabelnya terbelah menjadi \emph{variabel poros} dan
        \emph{variabel \omterm{def:b1:vspaces:free}{bebas}} (yaitu parameternya); lalu penyulihan mundur
        menyatakan yang pertama lewat yang kedua: sehingga \omterm{def:b1:logic:sets}{himpunan} penyelesaiannya berupa
        \omterm{def:b1:vspaces:subspace}{subruang} afin berdimensi $=$ cacah variabel bebasnya.
\end{enumerate}
Rumus Cramer diperuntukkan bagi teori dan sistem kecil; sedangkan penghapusannya
merupakan algoritma yang praktis.
\end{method}

\begin{example}[Sebuah pembahasan dengan parameter]\label{ex:b1:det:parameter}
Untuk $m \in \R$, pandanglah
\[
\begin{cases}
x + y + mz = 1\\
x + my + z = 1\\
mx + y + z = 1 .
\end{cases}
\]
Matriksnya berdeterminan $-(m+2)(m-1)^2$ (yang dihitung pada
\cref{exo:b1:det:7} dengan menambahkan semua kolomnya pada yang pertama). Untuk $m
\neq 1, -2$:
penyelesaiannya tunggal $x = y = z = \frac{1}{m+2}$ (lewat kesetangkupannya). Untuk $m =
1$: satu persamaan yang berulang tiga kali, jadi sebuah bidang penyelesaian. Untuk $m =
-2$: menambahkan ketiga persamaannya memberikan $0 = 3$, jadi tak ada penyelesaian.
\end{example}

\begin{remark}[Jebakan yang lazim]
\label{rem:b1:det:pitfalls}
\emph{Determinannya tak \omterm{def:b1:linmaps:def}{linear} pada matriksnya}: $\det(A + B)
\neq \det A + \det B$ (sudah $\det(I_2 + I_2) = 4 \neq 2$); melainkan ia
\omterm{def:b1:linmaps:def}{linear} pada setiap \emph{kolomnya} secara terpisah, yang merupakan hal yang sama sekali
berbeda.
\emph{Penskalaan}: $\det(\lambda A) = \lambda^n\det A$, bukan
$\lambda\det A$ --- karena masing-masing dari $n$ kolomnya terskalakan.
\emph{\omterm{met:b1:matrices:gauss}{Operasi baris} tak semuanya \omterm{def:b1:vspaces:free}{bebas}}: $L_i \leftarrow L_i +
\lambda L_j$ mengawetkan determinannya, tetapi sebuah penukaran mengubah
tandanya dan $L_i \leftarrow \lambda L_i$ mengalikannya dengan $\lambda$
--- adapun galat pembukuan di sini merupakan sumber klasik tanda
yang salah pada perhitungan berbasis penghapusan.
\emph{Determinan nol itu awalnya, bukan akhirnya}: karena ia mengatakan
``rank $< n$'' tetapi bukan rank yang mana; hanya kerja lanjutan (yaitu bentuk
eselon, atau minor pada \cref{exo:b1:det:12}) yang menemukannya --- bandingkanlah
kasus $m = 1$ dengan $m = -2$ pada \cref{ex:b1:det:parameter}.
\emph{Cramer menuntut keterbalikan}: karena ketika $\det A = 0$
rumus $x_j = \det A_j/\det A$ tak bermakna, dan sistemnya
boleh jadi tetap berpenyelesaian (bahkan tak berhingga banyaknya).
\emph{Hanya matriks persegi yang berdeterminan}: bagi sistem yang
persegi panjang, penghapusannya satu-satunya alat.
\end{remark}

\begin{remark}[Ke mana determinannya pergi]
\label{rem:b1:det:whereused}
Tiga kehidupan menanti skalar ini. \emph{Geometri}: $\abs{\det}$ adalah
faktor penskalaan luas atau volume \omterm{def:b1:logic:map}{pemetaan} yang berkaitan --- yang dipertepat
bagi bidang pada \cref{ch:b1:euclid} dan, sebagai Jacobian
sebuah perubahan variabel, pada \omterm{thm:b1:integration:def}{integral} rangkap jilid
Tahun~2. \emph{Aljabar}: $\det(A - \lambda I)$, yaitu
\omterm{def:b1:poly:def}{polinomial} karakteristiknya, membuka teori nilai eigen pada Tahun~2 ---
dengan kesamaan $A^2 - (\operatorname{tr} A)A + (\det A)I = 0$ pada
soal akhir pekan \cref{ch:b1:matrices} sebagai bayangan pertamanya.
\emph{Analisis}: determinan matriks yang istimewa (Vandermonde,
Cauchy, Gram) memutuskan kapan masalah interpolasi, penguraian dan
\omterm{def:b1:linmaps:projection}{proyeksi} terajukan dengan baik; adapun soal akhir pekan di bawah
menilai kedua keluarga yang pertama selengkapnya.
\end{remark}

\begin{remark}[Cakrawala di dalam Buku 3]
\label{rem:b1:det:perspectives}
Bab ini menutup tulang punggung aljabar \omterm{def:b1:linmaps:def}{linear} jilid ini, dan
kedua bab sisanya mencairkan dividennya. Pada
\cref{ch:b1:euclid}: matriks Gram
$\bigl(\langle v_i, v_j\rangle\bigr)$ menguji kebebasan lewat sebuah
determinan (\cref{exo:b1:euclid:11}), dan isometri bidangnya
terbelah menjadi rotasi dan pencerminan menurut tanda
determinannya --- dan penggolongan soal akhir pekan di sana
berjalan di atasnya. Pada \cref{ch:b1:multivar}: besaran Monge $rt -
s^2$ merupakan determinan matriks setangkup \omterm{def:b1:derivative:def}{turunan}
keduanya, dan persamaan normal kuadrat terkecilnya berupa
sistem Cramer yang matriksnya matriks Gram (jadi matriks momen) ---
yang dapat dibalik justru lewat kriteria bercita rasa Vandermonde
yang ditegakkan di sini. Ketika bab itu menegaskan ``dapat dibalik'' atau
``positif'', kuitansinya ada pada bab ini.
\end{remark}

\section{Latihan}

\begin{exercise}[$\star$]\label{exo:b1:det:1}
Hitunglah:
\[
\begin{vmatrix} 3 & 1\\ 5 & 2 \end{vmatrix},
\qquad
\begin{vmatrix} 1 & 2 & 3\\ 4 & 5 & 6\\ 7 & 8 & 9\end{vmatrix},
\qquad
\begin{vmatrix} 1 & 1 & 1\\ 1 & 2 & 4\\ 1 & 3 & 9\end{vmatrix}.
\]
\end{exercise}

\begin{exercise}[$\star$]\label{exo:b1:det:2}
Untuk $\lambda \in \R$ yang mana keluarga $\bigl((1, 1, \lambda),
(1, \lambda, 1), (\lambda, 1, 1)\bigr)$ merupakan \omterm{def:b1:vspaces:free}{basis} $\R^3$?
\end{exercise}

\begin{exercise}[$\star$]\label{exo:b1:det:3}
Pecahkanlah lewat kaidah Cramer:
\[
\begin{cases}
2x + y = 5\\
3x - 2y = 4 ,
\end{cases}
\qquad\text{lalu}\qquad
\begin{cases}
x + y + z = 6\\
x - y + z = 2\\
2x + y - z = 1 .
\end{cases}
\]
\end{exercise}

\begin{exercise}[$\star$]\label{exo:b1:det:4}
Pecahkanlah lewat \omterm{met:b1:det:gauss}{penghapusan Gauss}, dengan memerikan \omterm{def:b1:logic:sets}{himpunan} penyelesaiannya:
\[
\begin{cases}
x + 2y - z + t = 1\\
2x + 4y + z - t = 5\\
x + 2y + 2z - 2t = 4 .
\end{cases}
\]
\end{exercise}

\begin{exercise}[$\star\star$]\label{exo:b1:det:5}
Buktikan rumus Vandermonde pada \cref{ex:b1:det:vandermonde} lewat
induksi pada $n$, dengan operasi kolom $C_k \leftarrow C_k -
x_1 C_{k-1}$ yang dijalankan dari $k = n$ turun ke $k = 2$.
\end{exercise}

\begin{exercise}[$\star\star$]\label{exo:b1:det:6}
(Tridiagonal) Misalkan $D_n$ determinan $n \times n$ dengan $2$ pada
diagonalnya, $1$ pada kedua diagonal yang bersebelahan, dan $0$ di tempat lain.
Dengan menjabarkannya sepanjang baris pertamanya, buktikanlah $D_n = 2D_{n-1} - D_{n-2}$ lalu
hitunglah $D_n$ ($D_1 = 2$, $D_2 = 3$).
\end{exercise}

\begin{exercise}[$\star\star$]\label{exo:b1:det:7}
Lengkapilah \cref{ex:b1:det:parameter}: hitunglah determinan
$\begin{vmatrix} 1 & 1 & m\\ 1 & m & 1\\ m & 1 & 1\end{vmatrix}$
lewat operasi $C_1 \leftarrow C_1 + C_2 + C_3$, lalu jalankanlah
pembahasan sistemnya selengkapnya.
\end{exercise}

\begin{exercise}[$\star\star$]\label{exo:b1:det:8}
Misalkan $A \in \mathcal{M}_n(\R)$ berentri \emph{bilangan bulat}. Buktikan
bahwa $A$ mempunyai invers berentri bilangan bulat jika dan hanya jika $\det A
= \pm 1$. \emph{(Untuk arah langsungnya, ambillah determinannya; sedangkan untuk
konversnya, akuilah --- atau buktikanlah untuk $n \leq 3$ lewat kofaktor --- bahwa
$A^{-1} = \frac{1}{\det A}\,\operatorname{Com}(A)^{\mathsf T}$ dengan
matriks kofaktor yang bulat.)}
\end{exercise}

\begin{exercise}[$\star\star\star$]\label{exo:b1:det:9}
Hitunglah determinan $n \times n$ atas matriks $aI + bJ$
(\cref{exo:b1:matrices:9}), yaitu dengan $a + b$ pada diagonalnya dan
$b$ di tempat lain. \emph{(Tambahkanlah semua kolomnya pada yang pertama, faktorkanlah, lalu
bersihkanlah.)} Pulihkanlah syarat keterbalikannya $a \neq 0$, $a + nb
\neq 0$.
\end{exercise}

\begin{exercise}[$\star\star\star$]\label{exo:b1:det:10}
Misalkan $A, B \in \mathcal{M}_n(\R)$. Buktikan bahwa
\[
\det\begin{pmatrix} A & B\\ B & A \end{pmatrix}
= \det(A + B)\,\det(A - B),
\]
lewat operasi kolom dan baris blok ($C_1 \leftarrow C_1 + C_2$, lalu
$L_2 \leftarrow L_2 - L_1$, dalam bentuk blok), dengan mengandaikan kaidah
blok segitiga yang alami $\det\begin{pmatrix} M & N\\ 0 &
P\end{pmatrix} = \det M \det P$ --- yang dibuktikan bagi blok $2 \times 2$
pada \cref{ex:b1:det:blocktriangular}.
\end{exercise}

\begin{exercise}[$\star\star$]\label{exo:b1:det:11}
(Sirkulan berorde $3$) Misalkan $a, b, c \in \C$ dan
\[
\Delta = \begin{vmatrix}
a & b & c\\
c & a & b\\
b & c & a
\end{vmatrix}.
\]
Buktikan bahwa $\Delta = (a + b + c)(a^2 + b^2 + c^2 - ab - bc -
ca)$, lalu faktorkanlah selengkapnya atas $\C$ dengan memakai $j =
\eu^{2\iu\pi/3}$:
\[
\Delta = (a + b + c)(a + jb + j^2c)(a + j^2b + jc) .
\]
\emph{(Mulailah dengan $C_1 \leftarrow C_1 + C_2 + C_3$; sedangkan untuk
bentuk kompleksnya, perhatikanlah bahwa kolom $(1, j, j^2)^{\mathsf T}$
nyaris berperilaku bagai vektor eigen.)}
\end{exercise}

\begin{exercise}[$\star\star\star$]\label{exo:b1:det:12}
(Rank dan minor) Misalkan $A \in \mathcal{M}_{n,p}(K)$. Buktikan bahwa
$\operatorname{rk} A$ sama dengan ukuran terbesar $r$ atas sebuah
submatriks $r \times r$ pada $A$ yang dapat dibalik (dengan submatriks menyimpan
entri di persilangan $r$ baris terpilih dan $r$ kolom
terpilih). \emph{(Jika $\operatorname{rk} A = r$, pilihlah $r$ kolom
\omterm{def:b1:vspaces:free}{bebas}, lalu $r$ baris \omterm{def:b1:vspaces:free}{bebas} pada blok $n \times r$ yang dihasilkannya;
sebaliknya, sebuah submatriks yang dapat dibalik memaksa kolom
$A$ yang bersesuaian menjadi \omterm{def:b1:vspaces:free}{bebas}.)}
\end{exercise}

\section{Soal: alternan ganda Cauchy}

\begin{problem}\label{pb:b1:det:1}
Dua determinan menguasai penerapan bab ini: yaitu
\omterm{ex:b1:det:vandermonde}{determinan Vandermonde}, yang dinilai pada \cref{exo:b1:det:5}, dan
\emph{\omterm{pb:b1:det:1}{determinan Cauchy}} $\det\bigl(\frac{1}{a_i +
b_j}\bigr)$, yang dinilai di sini. Di sekitar keduanya soal ini menghimpun
kotak perkakas alternannya: yaitu kiat kolom \omterm{def:b1:poly:def}{polinomial}, interpolasi
lewat Cramer, \omterm{pb:b1:det:1}{matriks Hilbert}, diskriminan sebuah kubik, dan
metode \omterm{def:b1:poly:def}{polinomial} berselang-seling. Sepanjang soal ini, $V(x_1,
\dots, x_n) = \prod_{i < j}(x_j - x_i)$ menyatakan nilai
Vandermondenya.\index{determinan Cauchy}\index{matriks Hilbert}
\index{diskriminan}

\medskip
\noindent\textbf{Bagian I --- Kotak perkakas Vandermonde.}
\begin{enumerate}
  \item Hitunglah $V(1, 2, 3, 4)$, lalu ingatlah kembali mengapa interpolasi pada
        $n$ simpul yang berbeda berpasangan merupakan sistem Cramer.
  \item (Alternan \omterm{def:b1:poly:def}{polinomial}) Misalkan $P_0, \dots, P_{n-1}$
        bersifat \emph{\omterm{def:b1:poly:def}{monik}} dengan $\deg P_k = k$. Buktikan
        \[
        \det\bigl(P_{j-1}(x_i)\bigr)_{1 \leq i, j \leq n}
        = V(x_1, \dots, x_n) :
        \]
        karena operasi kolom mengganti setiap kolom pangkatnya dengan sebarang tangga
        \omterm{def:b1:poly:def}{monik}, secara cuma-cuma.
  \item Terapkanlah pertanyaan 2 pada \omterm{def:b1:poly:def}{polinomial} binomial $B_k =
        \frac{X(X-1)\cdots(X-k+1)}{k!}$: buktikanlah bahwa untuk
        \emph{bilangan bulat} $m_1 < m_2 < \dots < m_n$,
        \[
        \frac{V(m_1, \dots, m_n)}{0!\,1!\,2!\cdots(n-1)!}
        \in \N :
        \]
        hasil kali semua selisih berpasangan atas $n$ bilangan bulat
        terbagi oleh superfaktorial $0!\,1!\cdots(n-1)!$.
  \item Buktikan $\det\bigl(x_i^{\,j}\bigr)_{1 \leq i, j \leq n} =
        x_1 x_2 \cdots x_n\, V(x_1, \dots, x_n)$ (dengan pangkatnya kini
        bermula di $1$).
  \item (Matriks momen) Misalkan $S = \bigl(p_{i+j-2}\bigr)_{1 \leq
        i, j \leq n}$ dengan $p_k = x_1^k + \dots + x_n^k$. Buktikan
        bahwa $S = W^{\mathsf T} W$ bagi matriks $W =
        (x_i^{\,j-1})_{ij}$, lalu simpulkanlah
        \[
        \det S = V(x_1, \dots, x_n)^2 ,
        \]
        lalu simpulkanlah: bahwa $n$ bilangan \emph{real} berbeda
        berpasangan jika dan hanya jika matriks momennya
        dapat dibalik, dan bahwa $\det S \geq 0$ selalu.
\end{enumerate}

\medskip
\noindent\textbf{Bagian II --- Interpolasi ditinjau ulang.} Dengan simpul
$x_1 < \dots < x_n$, dan nilai $y_1, \dots, y_n$.
\begin{enumerate}[resume]
  \item Tulislah syarat ``$P = c_0 + c_1X + \dots +
        c_{n-1}X^{n-1}$ menginterpolasi'' sebagai \omterm{def:b1:det:system}{sistem linear} pada
        $c_k$ dengan matriks $W$, lalu pulihkanlah dari $\det W = V \neq
        0$ keberadaan dan ketunggalan interpolannya
        (bandingkanlah kedua bukti yang terdahulu,
        \cref{thm:b1:poly:lagrange} dan
        \cref{ex:b1:linmaps:interpolation}).
  \item Lewat kaidah Cramer dan \omterm{thm:b1:det:cofactor}{penjabaran kofaktor} atas determinan
        yang bersangkutan sepanjang kolom terakhirnya, buktikanlah bahwa
        koefisien utama interpolannya adalah
        \[
        c_{n-1} = \sum_{i=1}^{n}
        \frac{y_i}{\prod_{j \neq i}(x_i - x_j)} .
        \]
  \item (Vandermonde yang menyatu) Hitunglah
        \[
        \begin{vmatrix}
        1 & x_1 & x_1^2\\
        0 & 1 & 2x_1\\
        1 & x_2 & x_2^2
        \end{vmatrix}
        = (x_2 - x_1)^2 ,
        \]
        lalu tafsirkanlah: bahwa data $\bigl(P(x_1), P'(x_1),
        P(x_2)\bigr)$ menentukan satu $P \in \R_2[X]$ yang tunggal ketika
        $x_1 \neq x_2$ (yaitu interpolasi Hermite).
  \item Carilah $P \in \R_2[X]$ yang tunggal dengan $P(0) = 1$, $P'(0)
        = 0$, $P(1) = 2$, lalu periksalah jawabanmu terhadap
        pertanyaan 8.
\end{enumerate}

\medskip
\noindent\textbf{Bagian III --- \omterm{pb:b1:det:1}{Determinan Cauchy}.} Misalkan $a_1,
\dots, a_n$ dan $b_1, \dots, b_n$ skalar dengan $a_i + b_j
\neq 0$ untuk semua $i, j$, dan
\[
C_n = \det\Bigl(\frac{1}{a_i + b_j}\Bigr)_{1 \leq i, j \leq n} .
\]
\begin{enumerate}[resume]
  \item Hitunglah $C_2$ dengan tangan lalu tempatkanlah ia dalam bentuk
        ``hasil kali selisih dibagi hasil kali jumlah''.
  \item Untuk $n \geq 2$, jalankanlah $L_i \leftarrow L_i - L_n$ ($i <
        n$) lalu faktorkanlah baris dan kolomnya untuk membuktikan
        \[
        C_n = \frac{\prod_{i<n}(a_n - a_i)}{\prod_{j}(a_n +
        b_j)}\;\det M,
        \]
        dengan $M$ yang sepakat dengan matriks Cauchy pada baris $i < n$
        dan berbaris terakhir $(1, 1, \dots, 1)$.
  \item Jalankanlah $C_j \leftarrow C_j - C_n$ ($j < n$) pada $M$,
        faktorkanlah lagi, lalu simpulkanlah lewat induksi
        \emph{alternan ganda Cauchy}:
        \[
        C_n = \frac{\prod_{1 \leq i < j \leq n}(a_j - a_i)(b_j -
        b_i)}{\prod_{i, j}(a_i + b_j)} .
        \]
  \item Simpulkanlah kriteria keterbalikannya (yaitu $a_i$ berbeda
        berpasangan dan $b_j$ berbeda berpasangan). Bagi
        \emph{\omterm{pb:b1:det:1}{matriks Hilbert}} $H_n = \bigl(\frac{1}{i + j -
        1}\bigr)$: hitunglah $\det H_2$ dan $\det H_3$ dari
        rumusnya, lalu periksalah bahwa $H_2^{-1}$ berentri bilangan bulat.
  \item Tunjukkan bahwa untuk $b_j$ yang berbeda berpasangan dan sebarang ruas
        kanan, sistem $\sum_j \frac{c_j}{a_i + b_j} = y_i$
        ($i = 1, \dots, n$) berpenyelesaian tunggal, lalu kaitkanlah
        ini dengan keberadaan dan ketunggalan penguraian pecahan
        parsial berkutub sederhana
        (\cref{thm:b1:fractions:complex}).
\end{enumerate}

\medskip
\noindent\textbf{Bagian IV --- Diskriminan sebuah kubik.} Misalkan
$\lambda_1, \lambda_2, \lambda_3$ akar (di dalam $\C$) atas $X^3
+ pX + q$, dan $p_k = \lambda_1^k + \lambda_2^k + \lambda_3^k$.
\begin{enumerate}[resume]
  \item Dengan memakai $\lambda^3 = -p\lambda - q$ pada setiap akarnya dan
        Vieta ($p_1 = 0$), hitunglah $p_2 = -2p$, $p_3 = -3q$, dan
        $p_4 = 2p^2$.
  \item Dengan pertanyaan 5 (atas $\C$, dengan menjaga $\det S = V^2$),
        hitunglah
        \[
        \operatorname{disc} = V(\lambda_1, \lambda_2,
        \lambda_3)^2 = \begin{vmatrix}
        3 & 0 & -2p\\
        0 & -2p & -3q\\
        -2p & -3q & 2p^2
        \end{vmatrix}
        = -4p^3 - 27q^2 .
        \]
  \item Simpulkanlah: bahwa $X^3 + pX + q$ berakar rangkap jika dan hanya
        jika $4p^3 + 27q^2 = 0$; lalu periksalah pada $X^3 - 3X + 2 = (X -
        1)^2(X + 2)$.
  \item Andaikan $p, q$ real. Buktikan bahwa kubiknya berakar real
        tiga yang berbeda jika dan hanya jika $\operatorname{disc}
        > 0$, dan berakar satu real tambah dua akar sekawan yang tak real jika
        dan hanya jika $\operatorname{disc} < 0$. \emph{(Jika
        $\lambda_3 = \conj{\lambda_2} \neq \lambda_2$ dan
        $\lambda_1 \in \R$, tunjukkanlah bahwa $V$ murni khayal.)}
\end{enumerate}

\medskip
\noindent\textbf{Bagian V --- Dividen, dan metode
berselang-selingnya.}
\begin{enumerate}[resume]
  \item Untuk $0 < a_1 < a_2 < \dots < a_n$, tunjukkanlah
        $\det\bigl(\frac{1}{a_i + a_j}\bigr) > 0$.
  \item Hitunglah $\det\bigl(\binom{m_i}{j-1}\bigr)_{1 \leq i, j
        \leq 3}$ untuk $(m_1, m_2, m_3) = (2, 4, 7)$, mula-mula lewat
        pertanyaan 2--3, lalu lewat penjabaran langsung.
  \item Misalkan $\lambda_1, \dots, \lambda_n$ berbeda berpasangan
        dan taknol. Dengan memakai matriks Vandermonde yang dapat dibalik,
        buktikanlah lagi bahwa barisan geometri
        $\bigl((\lambda_i^{\,k})_{k \geq 0}\bigr)_{1 \leq i \leq
        n}$ membentuk \omterm{def:b1:vspaces:free}{keluarga bebas} pada ruang barisannya.
  \item Hitunglah $\det\bigl(\frac{1}{i + j}\bigr)_{1 \leq i, j
        \leq 3}$ dari alternan gandanya.
  \item (\omterm{def:b1:poly:def}{Polinomial} berselang-seling) Sebutlah sebuah \omterm{def:b1:poly:def}{polinomial} $F$ pada $x_1,
        \dots, x_n$ \emph{berselang-seling} bila menukar dua
        variabelnya mengubah tandanya. Tunjukkan bahwa $F$ yang berselang-seling
        lenyap kapan pun $x_i = x_j$ ($i \neq j$), lalu simpulkanlah
        --- satu variabel setiap kali, lewat teorema faktornya ---
        bahwa $F$ terbagi oleh $\prod_{i<j}(x_j - x_i)$.
  \item Pakailah pertanyaan 23 untuk membuktikan ulang rumus Vandermonde tanpa
        induksi: bahwa determinan $\det(x_i^{\,j-1})$ merupakan
        \omterm{def:b1:poly:def}{polinomial} berselang-seling berderajat total $\binom n2$,
        jadi kelipatan \emph{tetap} atas $\prod_{i<j}(x_j -
        x_i)$; lalu kenalilah tetapannya dengan membandingkan satu monomial.
  \item Sintesis, dalam empat kalimat: sifat tunggal yang mana pada
        determinannya (yaitu aksioma yang mana) yang menghasilkan semua
        pemfaktoran soal ini; mengapakah kesamaan matriks momen
        pada pertanyaan 5 mengubah \omterm{def:b1:logic:statement}{pernyataan} tentang
        keberbedaan \emph{kompleks} menjadi uji tanda
        \emph{real} yang dapat dihitung; kedua matriks klasik yang mana yang
        dinilai selengkapnya di sini dan masalah \omterm{def:b1:linmaps:def}{linear} yang mana yang
        dikuasainya; dan bagaimanakah metode berselang-seling pertanyaan
        23--24 menjelaskan, dalam satu gerakan, mengapa $\prod_{i<j}(x_j -
        x_i)$ terus muncul. Namailah teorema Bagian III.

\end{enumerate}
\end{problem}
